% E100_arithmetic_ground_truth.tex
%
% Stand-alone proposition: the arithmetic ground truth for the elliptic curve
% E_{100/Q} = LMFDB 100.a1 at the five quintic falsifier primes.
% Used by chapters/examples/cy_c_six_routes_convergence.tex
% (sec:cy-c-class-B-quintic, rem:quintic-hecke-falsifier).
%
% This is a stand-alone arithmetic certificate; it is NOT used by main.tex.
% It is picked up by compute/scripts/audit_independent_verification.py.
%
% Author: Raeez Lorgat. 2026-04-17.

\documentclass{article}
\usepackage{amsmath, amssymb, amsthm, enumitem}
\newtheorem{proposition}{Proposition}
\providecommand{\ClaimStatusProvedHere}{}

\begin{document}

\begin{proposition}[$E_{100/\mathbb{Q}}$ arithmetic ground truth at the
quintic falsifier primes]
\label{prop:E100-arithmetic-ground-truth}
\ClaimStatusProvedHere
The Hecke eigenvalues of the elliptic curve $E_{100/\mathbb{Q}} =
\textup{LMFDB } \texttt{100.a1}$ at the five quintic falsifier primes
$\{3,\, 7,\, 13,\, 29,\, 37\}$ are
\[
  a_3 = 2,\quad a_7 = -2,\quad a_{13} = -2,\quad a_{29} = 6,\quad a_{37} = -2,
\]
each obtained as $a_p = p + 1 - \#E(\mathbb{F}_p)$ from the minimal
Weierstrass model $[a_1, a_2, a_3, a_4, a_6] = [0,\, -1,\, 0,\, -33,\, 62]$.
The values satisfy:
\begin{enumerate}[label=\textup{(\roman*)}]
\item the Hasse bound $|a_p| \leq 2\sqrt{p}$ at each of the five primes;
\item Hecke multiplicativity $a_{mn} = a_m \cdot a_n$ for every coprime
      pair $(m, n)$ tested with $\max(m, n) \leq 100$;
\item the prime-power recursion $a_{p^2} = a_p^{\,2} - p$ for every good
      prime $p \leq 50$;
\item rank-$0$ analytic positivity: the Mellin-truncated partial sum
      $L(1,\, E_{100}) \approx 0.6315 > 0$, consistent with the
      LMFDB-listed analytic rank $0$ for \texttt{100.a1};
\item discriminant--conductor compatibility: $\Delta = 1\,250\,000 =
      2^4 \cdot 5^7$ supports exactly the bad primes $\{2, 5\}$ of
      conductor $N = 100 = 2^2 \cdot 5^2$.
\end{enumerate}
\end{proposition}

\begin{proof}
Direct exhaustive enumeration of $(x, y) \in \mathbb{F}_p^2$ satisfying
$y^2 = x^3 - x^2 - 33 x + 62$ gives
\[
  \#E(\mathbb{F}_3) = 2,\quad \#E(\mathbb{F}_7) = 10,\quad
  \#E(\mathbb{F}_{13}) = 16,\quad \#E(\mathbb{F}_{29}) = 24,\quad
  \#E(\mathbb{F}_{37}) = 40,
\]
yielding $a_p = p + 1 - \#E(\mathbb{F}_p)$ as $\{2, -2, -2, 6, -2\}$.
The five corroborating constraints (i)--(v) are independent classical
theorems about $a_p$ that hold without re-deriving $a_p$ from the
Weierstrass equation:
\begin{itemize}[leftmargin=2em]
\item (i) is Hasse 1936 (\emph{Zur Theorie der abstrakten elliptischen
      Funktionenkörper}, J.~Reine Angew.~Math.~175), bounding
      $|a_p| \leq 2\sqrt{p}$ from the Riemann hypothesis for $E/\mathbb{F}_p$.
\item (ii)--(iii) are Hecke 1937 (\emph{Über die Bestimmung
      Dirichletscher Reihen}, Math.~Ann.~112): the Euler product of an
      eigenform L-series gives $a_{mn} = a_m a_n$ for $\gcd(m, n) = 1$
      and the prime-power recursion at every good prime.
\item (iv) is Coates--Wiles 1977 (Invent.~Math.~39) for the CM case
      together with the general modularity theorem (Wiles 1995,
      Ann.~Math.~141, completed by Breuil--Conrad--Diamond--Taylor 2001):
      $L(1, E) \neq 0$ iff $\mathrm{rank}\,E(\mathbb{Q}) = 0$. The
      partial sum is computed from the $a_p$ table via the
      Mellin-truncated formula
      $L(1, E) \approx \sum_n a_n / n \cdot \exp(-2\pi n / \sqrt{N})$
      with $N = 100$.
\item (v) is the N\'{e}ron--Ogg--Shafarevich criterion (N\'{e}ron 1964,
      Ogg 1967, Shafarevich): the conductor's bad-prime support is
      contained in the discriminant's bad-prime support. Equality forces
      no extra ramification.
\end{itemize}
The point-counted values are verified to satisfy each constraint,
jointly establishing the proposition. The engine
\texttt{compute/quintic\_E100\_falsifier.py} performs the computation
and the test suite
\texttt{compute/tests/test\_quintic\_E100\_falsifier.py} certifies
each independent constraint.
\end{proof}

\end{document}
